\begin{lemma}[Elementary Gaussian moment bounds on $\mathbb{R}^n$]
Fix $n \geq 1$ and work on $\mathbb{R}^n$ with Lebesgue measure.  Then all of the
following hold.
\begin{enumerate}
  \item For every $b > 0$ the function $w \mapsto e^{-b|w|^2}$ is Lebesgue integrable
    on $\mathbb{R}^n$.
  \item For every $b > 0$,
    \[
      \int_{\mathbb{R}^n} e^{-b|w|^2}\,dw = \Bigl(\frac{\pi}{b}\Bigr)^{n/2}.
    \]
  \item For every $w \in \mathbb{R}^n$ and every coordinate $j$, $|w_j| \leq |w|$.
  \item For every $b > 0$ and every $v \in \mathbb{R}^n$,
    \[
      |v|\,e^{-b|v|^2} \leq \Bigl(1 + \frac{2}{b}\Bigr)\,e^{-(b/2)|v|^2}.
    \]
  \item For every $b > 0$ the function $w \mapsto |w|\,e^{-b|w|^2}$ is Lebesgue
    integrable on $\mathbb{R}^n$, and
    \[
      \int_{\mathbb{R}^n} |w|\,e^{-b|w|^2}\,dw
        \leq \Bigl(1 + \frac{2}{b}\Bigr)\Bigl(\frac{\pi}{b/2}\Bigr)^{n/2}.
    \]
  \item For all $\sigma > 0$ and $b' > 0$,
    \[
      \pi^{-n/2}\,\sigma^{-n/2}\,\Bigl(\frac{\pi}{b'}\Bigr)^{n/2}
        = \Bigl(\frac{1}{\sigma b'}\Bigr)^{n/2}.
    \]
\end{enumerate}
\end{lemma}
\begin{proof}
Each item is proved separately; items (4) and (5) use the earlier ones, and no item
uses any later one.

\emph{(1).}  Fix $b > 0$.  The Gaussian $w \mapsto e^{-b|w|^2}$ on the $n$-dimensional
inner product space $\mathbb{R}^n$ is integrable: this is the real part of the standard
complex Gaussian integrability statement for $w \mapsto e^{-b|w + c|^2}$ with complex
parameter $b$ of positive real part, specialised to $c = 0$ and to real $b > 0$.  Since
$e^{-b|w|^2}$ is real valued, taking real parts changes nothing, and integrability of
the complex function gives integrability of this one.

\emph{(2).}  Fix $b > 0$.  The standard Gaussian integral over a real inner product
space of finite dimension $d$ is
$\int e^{-b|w|^2}\,dw = (\pi/b)^{d/2}$.  Here the space is $\mathbb{R}^n$, whose
dimension is $d = n$, which gives $(\pi/b)^{n/2}$.

\emph{(3).}  Fix $w$ and $j$.  Writing the Euclidean norm as
$|w| = \bigl(\sum_i |w_i|^2\bigr)^{1/2}$, the single term $|w_j|^2$ is at most the full
sum $\sum_i |w_i|^2$, because every term is nonnegative.  Taking square roots, which is
monotone, gives
\[
  |w_j| = \bigl(|w_j|^2\bigr)^{1/2} \leq \Bigl(\sum_i |w_i|^2\Bigr)^{1/2} = |w| .
\]

\emph{(4).}  Fix $b > 0$ and $v$, and set $t = |v| \geq 0$ and
$E = e^{-(b/2)t^2}$.  Then $E > 0$, and $E \leq 1$ because $(b/2)t^2 \geq 0$.  Also
$e^{-bt^2} = E^2$, since $-bt^2 = -(b/2)t^2 + (-(b/2)t^2)$.  Next, from the elementary
inequality $1 + u \leq e^u$ applied at $u = (b/2)t^2$ we get
$(b/2)t^2 \leq e^{(b/2)t^2}$, and since $e^{(b/2)t^2}E = 1$ this yields
\[
  \bigl((b/2)t^2\bigr)E \leq 1 .
\]
Multiplying by $2/b > 0$ and simplifying gives $t^2 E \leq 2/b$.  Now estimate, using
$t \leq 1 + t^2$ (equivalently $(t-1)^2 \geq 0$ together with $t \ge 0$) and $E^2 > 0$:
\[
  t\,e^{-bt^2} = t\,E^2 \leq (1 + t^2)E^2 = E^2 + (t^2E)E
    \leq 1 \cdot E + \frac{2}{b}\,E = \Bigl(1 + \frac{2}{b}\Bigr)E ,
\]
where the last inequality used $E \leq 1$ on the first summand and $t^2E \leq 2/b$ on
the second.  This is the claim, since $E = e^{-(b/2)|v|^2}$.

\emph{(5).}  Fix $b > 0$ and put $b/2 > 0$.  By (1) the function
$w \mapsto e^{-(b/2)|w|^2}$ is integrable, hence so is its constant multiple
$w \mapsto (1 + 2/b)e^{-(b/2)|w|^2}$; call that majorant $M$.  The function
$w \mapsto |w|e^{-b|w|^2}$ is continuous, hence strongly measurable, and by (4) it is
pointwise dominated in absolute value by $M$ (the function is nonnegative, so its
absolute value is itself).  Domination by an integrable majorant gives integrability,
which is the first half of the claim.  For the second half, monotonicity of the integral
applied to the pointwise bound of (4), followed by (2) at rate $b/2$, gives
\[
  \int |w|e^{-b|w|^2}\,dw
    \leq \int \Bigl(1 + \frac{2}{b}\Bigr)e^{-(b/2)|w|^2}\,dw
    = \Bigl(1 + \frac{2}{b}\Bigr)\int e^{-(b/2)|w|^2}\,dw
    = \Bigl(1 + \frac{2}{b}\Bigr)\Bigl(\frac{\pi}{b/2}\Bigr)^{n/2} .
\]

\emph{(6).}  Fix $\sigma, b' > 0$.  All bases below are positive, so the real power
laws apply.  Rewriting $-n/2 = -(n/2)$ and using
$c^{-n/2} = (c^{n/2})^{-1}$ for $c > 0$, together with
$(\pi/b')^{n/2} = \pi^{n/2}/(b')^{n/2}$ and
$(1/(\sigma b'))^{n/2} = \bigl((\sigma b')^{n/2}\bigr)^{-1}
  = \bigl(\sigma^{n/2}(b')^{n/2}\bigr)^{-1}$, both sides become quotients of products
of the positive quantities $\pi^{n/2}$, $\sigma^{n/2}$ and $(b')^{n/2}$.  The left side
is $\pi^{n/2}/\bigl(\pi^{n/2}\sigma^{n/2}(b')^{n/2}\bigr)$ and the right side is
$1/\bigl(\sigma^{n/2}(b')^{n/2}\bigr)$; these agree because $\pi^{n/2} \neq 0$.
\end{proof}
